% Part: computability
% Chapter: computability-theory
% Section: k-1

\documentclass[../../../include/open-logic-section]{subfiles}

\begin{document}

\olfileid[fr]{cmp}{thy}{k1}
\olsection{Un exemple de réduction}

Appliquons la \olref[ppr]{prop:reduce}.

\begin{prop}
\ollabel{prop:k1}
Soit
\[
K_1 = \Setabs{e}{\cfind{e}(0) \fdefined}.
\]
Alors $K_1$ est effectivement énumérable,
mais non calculable.
\end{prop}

\begin{proof}
Puisque
$K_1=\Setabs{e}{\lexists[s][T(e,0,s)]}$,
$K_1$ est effectivement énumérable par le
\olref[eqc]{thm:exists-char}.

Pour montrer que $K_1$ n'est pas calculable,
réduisons-y $K_0$.

\begin{explain}
Cette réduction demande une petite astuce :
les questions portant sur~$K_1$ concernent
seulement les calculs lancés sur l'entrée~$0$.
Supposons qu'un interlocuteur puisse répondre
à de telles questions; on appelle « oracle »
un moyen d'obtenir ces réponses. Pour savoir
si $\tuple{e,x}\in K_0$, c'est-à-dire si la
machine d'indice~$e$ s'arrête sur l'entrée~$x$,
construisons une autre machine, d'indice~$e_x$,
qui ignore toute entrée qu'on lui fournit et
exécute à la place la machine~$e$ sur~$x$.
Alors $e$ s'arrête sur~$x$ si et seulement si
$e_x$ s'arrête sur~$0$ (ou sur n'importe
quelle autre entrée). La réponse de l'oracle
pour~$e_x$ fournit donc la réponse à la
question initiale. C'est la réduction cherchée.
\end{explain}

À l'aide de la fonction universelle partielle
calculable, définissons la fonction ternaire~$f$
par
\[
f(x,y,z) \simeq \cfind{x}(y).
\]
Elle ignore entièrement son troisième argument.
Choisissons un indice~$e$ tel que
$f\simeq\cfind{e}[3]$; on a donc
\[
\cfind{e}[3](x,y,z) \simeq \cfind{x}(y).
\]
Par le théorème $s$-$m$-$n$, il existe une
fonction récursive primitive $s(e,x,y)$ telle que,
pour tout~$z$,
\begin{align*}
\cfind{s(e,x,y)}(z) & \simeq \cfind{e}[3](x,y,z) \\
& \simeq \cfind{x}(y).
\end{align*}

\begin{explain}
Dans le raisonnement intuitif, $s(e,x,y)$
est l'indice d'une machine qui ignore son
entrée~$z$ et calcule $\cfind{x}(y)$.
\end{explain}

En particulier,
\[
\cfind{s(e,x,y)}(0) \fdefined \quad \text{si et seulement si} \quad
\cfind{x}(y) \fdefined.
\]
Ainsi, $\tuple{x,y}\in K_0$ si et seulement si
$s(e,x,y)\in K_1$. Pour les codes de paires,
la réduction prend donc la valeur
\[
g(w) = s(e,(w)_0,(w)_1)
\]
et, pour les entiers qui ne codent pas une paire,
elle prend un indice fixé d'une fonction nulle
part définie. La reconnaissance des codes de
paires est calculable : on teste $\fn{SeqCode}(w)$ et $\len{w}=2$; la fonction totale ainsi
définie réduit bien $K_0$ à~$K_1$.\footnote{Complément éditorial : la formule source pour $g$ ne traite pas les entiers non codants. Nous ajoutons ce cas pour obtenir une réduction totale. Nous employons également $\simeq$ pour comparer les fonctions et valeurs partielles.}
\end{proof}

\end{document}
